\par\Needspace{8\baselineskip}
\originalchapter{简并压力与恒星结构}\label{ch:stars}
\originalsection{引力与压力的自洽平衡}
正常恒星的引力压缩先加热气体, 核反应随后补充能量, 气体与辐射压力共同抵抗收缩. 白矮星中核聚变停止, 主要由电子简并压力支撑; 中子星中需要中子简并和核相互作用, 更高致密度还要广义相对论. 照片展示对象外观, 真正的结构计算则需密度与压强自洽满足$\dd P/\dd r=-GM(r)\rho(r)/r^2$、$\dd M/\dd r=4\pi r^2\rho(r)$及状态方程. 本章先用均匀球说明尺度, 不把它当成精确恒星模型.

质量$M$, 半径$R$的均匀球引力能为$E_g=-3GM^2/(5R)$. 推导时将外壳$\dd M=4\pi r^2\rho\dd r$从无穷远搬到已积累质量$M(r)$周围, 积分$-\int GM(r)\dd M/r$即可. 固定组成的电子数$N_e=M/(\mu_e m_p)$, $\mu_e$是每个电子对应的平均核子数, 与电子磁矩记号在本章分开使用.

\originalsection{非相对论白矮星}
均匀电子数密度$n_e=N_e/(4\pi R^3/3)$, 零温动能$E_k=(3/5)N_e\epsilon_F$. 展开给
\[
E_k=\frac{A M^{5/3}}{R^2},\quad
A=\frac{3\hbar^2}{10m_e}(9\pi/4)^{2/3}(\mu_em_p)^{-5/3},
\quad E_{tot}=\frac{AM^{5/3}}{R^2}-\frac{3GM^2}{5R}.
\]
定质量对$R$求极值, 得$R_0=(10A/3G)M^{-1/3}$, 二阶导数在此为正. 增加质量反使半径缩小, 因为引力增强要求更高简并密度, 不像普通物体简单堆叠体积. 本模型没有最大稳定质量, 但它的前提是$p_F\ll m_ec$, 即密度不能无限高. 核动能相对电子动能小, 主要质量由核承担, 也是采用这一模型的条件.

对${}^4$He组成$\mu_e=2$, 可等价写$MR_0^3=2(9\pi)^2\hbar^6/(G^3m_e^3m_\alpha^5)$. 原课用Sirius B比较时, 均匀模型预测与观测半径同数量级, 更细的非相对论模型并不一定更接近; 这提示了遗漏相对论, 不能把一次拟合的接近程度当作模型正确的证明.

\originalsection{相对论简并与质量极限}
当$p_F\gtrsim m_ec$时, 用超相对论$\epsilon=cp=\hbar ck$. 同样状态计数给$D(\epsilon)=V\epsilon^2/(\pi^2\hbar^3c^3)$, $U=3N\epsilon_F/4$, $P=n\epsilon_F/4=U/(3V)$, 所以$P\propto n^{4/3}$而非$n^{5/3}$. 压力随密度增大较慢, 支撑机制变弱.

均匀球此时$E_k=A_RM^{4/3}/R$, $A_R=(3/4)\hbar c(9\pi/4)^{1/3}(\mu_em_p)^{-4/3}$. 总能量为$[A_RM^{4/3}-3GM^2/5]/R$. 系数为正时超相对论近似希望扩大$R$, 为负时希望坍缩, 都没有有限$R$处的极小. 分界质量尺度是
\[
M_*=\left(\frac{5A_R}{3G}\right)^{3/2}
\propto\frac{(\hbar c/G)^{3/2}}{(\mu_em_p)^2}.
\]
均匀近似只能给尺度和近似系数; 由自洽非均匀结构方程得到的Chandrasekhar极限约为$5.83M_\odot/\mu_e^2$, 对$\mu_e\simeq2$约$1.46M_\odot$. 该数值是课内引用的更精确模型结果, 本稿没有把均匀球推导冒充它的精确证明.

若电子与质子转化为中子, 改用中子质量的简并状态方程会使半径尺度降至十几公里. 核力与广义相对论对中子星最大质量重要, 原课给的约$2M_\odot$只是定性范围示例; 本册不将它作为对所有核物质状态方程都成立的常数. 普通恒星、白矮星、中子星和黑洞的区别在于支撑或坍缩机制, 不能只按表面温度排序.
